\label{ch:vacuum}

\section{Derivation of the action scale from the dodecaphasic state}

The action decade is determined intrinsically, without adopting it as an
inherited unit. The local determinantal functional is applied to the previously
constructed HMT output \(\alpha_{\HMT}\) and defines
\begin{equation}
\Sigma_H=\varphi_{\HMT}\alpha_{\HMT}^{16}.
\label{eq:action-determinantal-reader}
\end{equation}
The exponent \(16\) is the order of the local cokernel of the signed closure; its decimal valuation fixes
\begin{equation}
\operatorname{ord}_{10}(\Sigma_H)=-34.
\label{eq:action-decade}
\end{equation}

The origin of the exponent can be verified within this volume. On the
local orbit of four leaves, the Hadamard block acts
\begin{equation}
H_4=
\begin{pmatrix}
1&1&1&1\\
1&1&-1&-1\\
1&-1&1&-1\\
1&-1&-1&1
\end{pmatrix}.
\label{eq:action-hadamard}
\end{equation}

\begin{lemma}[Local index of the action line]
The form normal of Smith of \(H_4\) is
\begin{equation}
\operatorname{SNF}(H_4)=\operatorname{diag}(1,2,2,4).
\label{eq:action-smith}
\end{equation}
Consequently,
\[
\operatorname{coker}(H_4)\simeq
\Z/2\Z\oplus\Z/2\Z\oplus\Z/4\Z,
\qquad
|\operatorname{coker}(H_4)|=16.
\]
\end{lemma}

\begin{proof}
The greatest common divisor of the entries is one; that of the nonzero
\(2\times2\) minors is two; that of the \(3\times3\) minors is four; and
\(|\det H_4|=16\). The successive quotients of these determinantal divisors
are \(1,2,2,4\), which prove \cref{eq:action-smith}. The product of the
invariants gives the order of the cokernel.
\end{proof}

The norm of the section \(\alpha_{\HMT}\) over those sixteen sheets
produces \(\alpha_{\HMT}^{16}\); the self-scaling \(\varphi_{\HMT}\) acts
only once on the basis, not once per sheet. Hence
\cref{eq:action-determinantal-reader}. The certified comparisons
\begin{equation}
\alpha_{\HMT}^{16}<10^{-34}
<\varphi_{\HMT}\alpha_{\HMT}^{16}<10^{-33}
\label{eq:action-decade-bounds}
\end{equation}
prove, without introducing the SI unit, that the decimal order is
exactly \(-34\). The terminal evaluation begins with
\begin{equation}
\Sigma_H
=1.046243838104756580184229208303089\ldots\times10^{-34}.
\label{eq:action-sigma-value}
\end{equation}
The substitution of the kernel by three copies with exponent \(16^3\), or the repeated application of \(\varphi\) on each sheet, would alter the genealogy. Both operations constitute falsification criteria rather than equivalent conventions. The causal chain that is preserved is therefore,
\[
\APP\to\TRIT\to\TPK\to U_{12}^{\rm sign}
\to\alpha_{\HMT}\to\Sigma_H
\to10^{-34}\U_S\to
\{\hbar_{\rm pre},\hbar_{\rm ret}\}.
\]
This volume does not rederive \(\alpha_{\HMT}\), but develops the metrological geometry
subsequent to the selection of the type and decade by
\(\Sigma_H\). The integer \(-34\) is neither a mantissa nor an adjustment to the International
System, but the terminal valuation of the local functional.

\section{Oriented sections of the action line}

The realizations before and after the return constitute two sections
of the same action line:

\begin{equation}
\hbar_{\rm pre}=H_5\,10^{-34}\U_S,\qquad
\hbar_{\rm ret}=\eta_{\rm ret}\,10^{-34}\U_S.
\label{eq:action-twoway}
\end{equation}

The common decade determines the line, while the oriented
information is expressed through

\begin{equation}
R_{\rm act}=\frac{\hbar_{\rm pre}}{\hbar_{\rm ret}}
=\frac{H_5}{\eta_{\rm ret}},
\qquad
\xi_{\rm act}=\frac12\log R_{\rm act}.
\label{eq:action-rapidity}
\end{equation}

The first quantity provides a multiplicative coordinate and the second
an additive coordinate of rapidity type. They do not represent two independent
Planck constants, but the two orientations of the same section
with respect to the return.

If \(h_a=2\pi\hbar_a\), \(a\in\{\mathrm{pre},\mathrm{ret}\}\), a quantity proportional to \(h_a^{1/2}\) changes as \(R_{\rm act}^{1/2}\), one proportional to \(h_a^{-1/2}\) changes as \(R_{\rm act}^{-1/2}\), and a quotient in which the action cancels is a bidirectional invariant.

